\documentclass[11pt]{article}
\usepackage[a4paper,margin=1in]{geometry}
\usepackage{fontspec}
\usepackage[boldfont=false]{xeCJK}
\setCJKmainfont{FandolSong-Regular.otf}
\usepackage{amsmath}
\usepackage{amssymb}
\usepackage{array}
\usepackage{booktabs}
\usepackage{graphicx}
\usepackage{adjustbox}
\usepackage{enumitem}
\setlist[itemize]{topsep=2pt,partopsep=0pt,parsep=0pt,itemsep=2pt,leftmargin=1.4em}
\setlist[enumerate]{topsep=2pt,partopsep=0pt,parsep=0pt,itemsep=2pt,leftmargin=1.6em}
\usepackage{parskip}
\usepackage{fvextra}
\fvset{breaklines=true,breakanywhere=true,fontsize=\small}
\setlength{\parindent}{0pt}

\begin{document}

\begin{center}
  {\LARGE\bfseries Characteristics and classification of living organisms}\\[0.35em]
  {\large Igcse Biology}
\end{center}
\vspace{0.6em}

\section*{The seven characteristics of living organisms}

Every living thing shares the same seven \textbf{characteristics} 特征 — seven \textbf{life processes} 生命过程 that all \textbf{organisms} 生物体 carry out. You can remember them with the letters \textbf{MRS GREN}. The examiner gives marks for the exact wording, so learn each definition.

\begin{itemize}
\item \textbf{Movement} 运动 — an action by an organism, or part of an organism, that changes its position or place.

\item \textbf{Respiration} 呼吸作用 — the chemical reactions in \textbf{cells} 细胞 that break down \textbf{nutrient molecules} 营养物质 and release \textbf{energy} 能量 for \textbf{metabolism} 新陈代谢.

\item \textbf{Sensitivity} 应激性 — the ability to detect and respond to changes in the internal or external \textbf{environment} 环境.

\item \textbf{Growth} 生长 — a permanent increase in size and \textbf{dry mass} 干重.

\item \textbf{Reproduction} 生殖 — the processes that make more of the same kind of organism.

\item \textbf{Excretion} 排泄 — the removal of the \textbf{waste products} 废物 of metabolism, and of substances the body has in excess (more than it needs).

\item \textbf{Nutrition} 营养 — the taking in of materials for energy, growth and development.

\end{itemize}

Metabolism means all the chemical reactions that happen inside an organism's cells.

\section*{Classification: sorting living things into groups}

There are millions of kinds of living thing. To study them, scientists \textbf{classify} 分类 them — they sort them into groups by the features they share. Organisms in the same group share more features with each other than with organisms in other groups.

\subsection*{Species and the binomial system}

The smallest common group is the \textbf{species} 物种. A species is a group of organisms that can reproduce to make \textbf{fertile offspring} 可育后代 ("fertile" means the offspring can themselves go on to have young).

Each species has a two-part \textbf{scientific name} 学名. This way of naming is the \textbf{binomial system} 双名法 — a system agreed all over the world, so every scientist uses the same name.

\begin{itemize}
\item The \textbf{first} part is the \textbf{genus} 属 (a larger group that the species belongs to). It starts with a capital letter.

\item The \textbf{second} part is the species. It starts with a small letter.

\end{itemize}

The whole name is written in italics, for example \emph{Homo sapiens} (humans). Two rats named \emph{Rattus norvegicus} and \emph{Rattus rattus} belong to the same genus but are different species.

\subsection*{Dichotomous keys}

A \textbf{dichotomous key} 二歧检索表 is a tool that helps you identify an unknown organism. "Dichotomous" means "splitting into two". At each step the key gives you \textbf{two} choices about a feature. You pick the one that fits your organism, and that choice sends you on to the next pair of choices. You repeat this until you reach the organism's name.

When you build a key, choose clear features you can see — for example "has wings / has no wings" — not features that change, such as size.

\subsection*{Classification and DNA (Supplement)}

Modern classification tries to show \textbf{evolutionary relationships} 进化关系 — how closely different organisms are related. Organisms that are more closely related share a more recent \textbf{ancestor} 祖先.

To measure this, scientists compare the \textbf{base sequence} 碱基序列 of DNA (the order of the chemical "letters" that make up an organism's \textbf{genes} 基因). The rule is simple:

\begin{itemize}
\item Organisms that share a \textbf{more recent} ancestor (more closely related) have \textbf{more similar} DNA base sequences.

\item Organisms that share only a \textbf{distant} ancestor have \textbf{less similar} base sequences.

\end{itemize}

So a DNA base sequence can be used both to classify organisms and to work out how closely two species are related.

\section*{Kingdoms}

The biggest groups that organisms are sorted into are called \textbf{kingdoms} 界. At Core level you place organisms into the first two kingdoms below; for Supplement you need all \textbf{five kingdoms}. You decide where an organism belongs by looking at its main features.

\begin{center}
\adjustbox{max width=\linewidth}{%
\begin{tabular}{ll}
\toprule
\textbf{Kingdom} & \textbf{Main features} \\
\midrule
\textbf{Animal} 动物 & made of many cells; cells have no \textbf{cell wall} 细胞壁; cannot make their own food, so they feed on other organisms; most can move their whole body \\
\textbf{Plant} 植物 & made of many cells; cells have a cell wall; make their own food by \textbf{photosynthesis} 光合作用, so many cells contain \textbf{chloroplasts} 叶绿体 \\
\textbf{Fungus} 真菌 & for example \textbf{moulds} 霉菌, \textbf{mushrooms} 蘑菇 and \textbf{yeast} 酵母; have cell walls; cannot photosynthesise; feed on dead or living matter \\
\textbf{Prokaryote} 原核生物 & for example \textbf{bacteria} 细菌; \textbf{single-celled} 单细胞; cells have no \textbf{nucleus} 细胞核 \\
\textbf{Protoctist} 原生生物 & for example \emph{Amoeba} and \textbf{algae} 藻类; usually single-celled; cells do have a nucleus \\
\bottomrule
\end{tabular}}
\end{center}

The fungus, prokaryote and protoctist kingdoms are Supplement only.

\subsection*{Grouping the animal kingdom}

The animal kingdom is split first into two: animals with a backbone and animals without one.

\textbf{Vertebrates} 脊椎动物 are animals with a \textbf{backbone} 脊柱. There are five main groups.

\begin{center}
\adjustbox{max width=\linewidth}{%
\begin{tabular}{ll}
\toprule
\textbf{Group} & \textbf{Main features} \\
\midrule
\textbf{Mammals} 哺乳动物 & have \textbf{fur} 毛发 or hair; feed their young on milk \\
\textbf{Birds} 鸟类 & have \textbf{feathers} 羽毛 and a \textbf{beak} 喙; lay eggs with hard shells \\
\textbf{Reptiles} 爬行动物 & have dry skin covered with \textbf{scales} 鳞片; lay eggs with leathery shells on land \\
\textbf{Amphibians} 两栖动物 & have moist skin; lay eggs in water; the young live in water \\
\textbf{Fish} 鱼类 & have wet scales and \textbf{fins} 鳍; breathe using \textbf{gills} 鳃; lay eggs in water \\
\bottomrule
\end{tabular}}
\end{center}

\textbf{Arthropods} 节肢动物 are animals with no backbone. They have a hard outer skeleton (an \textbf{exoskeleton} 外骨骼) and legs that bend at \textbf{joints} 关节. There are four main groups.

\begin{center}
\adjustbox{max width=\linewidth}{%
\begin{tabular}{ll}
\toprule
\textbf{Group} & \textbf{Main features} \\
\midrule
\textbf{Insects} 昆虫 & body in 3 parts; 3 pairs of legs; usually 2 pairs of wings; 1 pair of \textbf{antennae} 触角 \\
\textbf{Arachnids} 蛛形类 & body in 2 parts; 4 pairs of legs; no wings; no antennae \\
\textbf{Crustaceans} 甲壳类 & many pairs of legs; 2 pairs of antennae; most live in water \\
\textbf{Myriapods} 多足类 & long body made of many \textbf{segments} 体节; one or two pairs of legs on each segment \\
\bottomrule
\end{tabular}}
\end{center}

\subsection*{Grouping the plant kingdom (Supplement)}

Plants are placed into groups too. You need two of them.

\begin{itemize}
\item \textbf{Ferns} 蕨类植物 — have roots, stems and leaves, but make \textbf{no} flowers or seeds. They reproduce using tiny \textbf{spores} 孢子.

\item \textbf{Flowering plants} 开花植物 — make flowers, and form \textbf{seeds} 种子 inside the flower. They split into two groups by their seed-leaves (the first leaves inside a seed, called \textbf{cotyledons} 子叶):

\end{itemize}

\begin{center}
\adjustbox{max width=\linewidth}{%
\begin{tabular}{ll}
\toprule
\textbf{Group} & \textbf{Features} \\
\midrule
\textbf{Monocotyledons} 单子叶植物 & seed has \textbf{one} cotyledon; long narrow leaves; \textbf{veins} 叶脉 run side by side (parallel) \\
\textbf{Dicotyledons} 双子叶植物 & seed has \textbf{two} cotyledons; broad leaves; veins form a branching net \\
\bottomrule
\end{tabular}}
\end{center}

\subsection*{Viruses (Supplement)}

\textbf{Viruses} 病毒 are not placed in any kingdom. They are not made of cells, so many scientists do not count them as living. A virus is very simple, with only two parts:

\begin{itemize}
\item an outer \textbf{protein coat} 蛋白质外壳, and

\item \textbf{genetic material} 遗传物质 (its genes) inside.

\end{itemize}

A virus cannot carry out the life processes on its own. It can only reproduce inside the living cells of a \textbf{host} 宿主 — the organism it infects.

\section*{Exam tips}

\begin{itemize}
\item Learn the \textbf{seven characteristics} by heart, with their exact definitions. Write them as clear, full sentences.

\item \textbf{Digestion} 消化 is \emph{not} one of the seven characteristics — it is just one part of nutrition. "Breathing" is not one of the seven either.

\item In a scientific name, the genus starts with a capital letter and the species with a small letter, and both are in italics.

\item To use a dichotomous key, start at the first pair of choices and follow the route step by step. Do not jump ahead.

\item More similar DNA base sequence → the organisms are more closely related (they share a more recent ancestor).

\end{itemize}


\end{document}
